\documentclass{lecturenotes}

\SetReferencesConfiguration{
  hyperref = true,
  hyperref/options = { unicode, hyperindex, pdfpagelabels },
  hyperref/setup = {
    breaklinks,
    colorlinks,
    linktocpage,
    urlcolor = blue,
    linkcolor = deepblue,
    filecolor = magenta,
    citecolor = blue,
    pdftitle = {Пример lecture notes},
    pdfauthor = {Desiment}
  },
  shortrefs = {
    picref = рис.,
    tabref = табл.,
    secref = разд.,
    lstref = листинг
  },
  zref-clever = true,
  zref-clever/setup = { lang = russian },
  azcref = true,
  biblatex = true,
  biblatex/options = { backend = biber, style = numeric, sorting = nyt, defernumbers = true },
  biblatex/path = .,
  biblatex/bibliographies = {
    { name = main, title = {Список литературы}, files = example.bib },
    { name = online, title = {Онлайн-ресурсы}, files = example-online.bib }
  },
  indexes = true,
  indexes/options = { makeindex, noautomatic },
  indexes/list = {
    { name = notation, title = {Список обозначений}, intoc, columns = 2 },
    { title = {Предметный указатель}, intoc, columns = 2 }
  }
}

\title{Демонстрационные лекционные заметки}
\author{Desiment}
\date{\today}

\begin{document}

\maketitle
\tableofcontents

\chapter{Инструменты для математического текста}

\section{Сюжет примера}
\label{sec:story}

Этот файл показывает, как класс \texttt{lecturenotes} собирает вместе
\texttt{tssuite}, \texttt{xamsmath}, \texttt{flsuite} и \texttt{refsuite}.
Короткие ссылки создаются конфигурацией \texttt{refsuite}: далее мы будем
ссылаться на \picref{fig:markov-chain}, \tabref{tab:operators},
\lstref{lst:recursion} и \secref{sec:references}.

\begin{definition}[note=Дискретная случайная величина,label=def:random-variable]
Случайная величина $X$ называется дискретной, если множество ее значений не
более чем счетно и для любого $A \subseteq \R$ вероятность события
$\set{X \in A}$ равна сумме вероятностей отдельных атомов.
\end{definition}

\begin{Theorem}[Формула полной вероятности][label=thm:total-probability]
Пусть события $H_1,\dots,H_n$ образуют разбиение пространства элементарных
исходов и $\P(H_i)>0$. Тогда для любого события $A$ верна формула
\[
  \P(A)=\sum_{i=1}^n \P(A | H_i)\P(H_i).
\]
\zlabel{zthm:total-probability}
\end{Theorem}

В силу \azcref{zthm:total-probability} вычисление сложной вероятности можно
разбить на несколько условных вычислений. Обычная ссылка на тот же результат:
теорема~\ref{thm:total-probability}.

\begin{Lemma}[Линейность математического ожидания][label=lm:expectation]
Для интегрируемых случайных величин $X$ и $Y$ и чисел $a,b\in\R$ выполнено
\[
  \E[aX+bY] = a\E[X] + b\E[Y].
\]
\end{Lemma}

\begin{remark}
Фреймированные окружения удобны для ключевых утверждений, а обычные окружения
подходят для промежуточных фактов и упражнений в тексте лекции.
\end{remark}

\section{Нотация xamsmath}
\label{sec:xamsmath}

Пакет \texttt{xamsmath} загружен классом со всеми основными плагинами. Он дает
единый стиль для множеств, алгебры, анализа, комбинаторики и вероятностей.

\index[notation]{$\N$}\index[notation]{$\R$}\index{математическая нотация}
\begin{align*}
  \set{x \in \R | x > 0} &\quad \text{множество положительных чисел}, &
  \card{A} &\quad \text{мощность множества}, \\
  \N,\Z,\Q,\R,\C &\quad \text{стандартные числовые множества}, &
  \norm{v}^2 &= \spr{v}{v}.
\end{align*}

Для анализа и вероятностей доступны компактные команды:
\begin{align*}
  \od{y}{x} &= 2x, &
  \pd{u}{t} &= \Delta u, \\
  \P(A | B) &= \frac{\P(A\cap B)}{\P(B)}, &
  \E[\Cov{X}{Y} | Z] &= 0.
\end{align*}

Комбинаторные обозначения остаются читаемыми даже в длинных формулах:
\[
  \sum_{k=0}^n \Stirling{n}{k}\risef{x}^k = x^n,
  \qquad
  \fallf{x}^{m+n}=\fallf{x}^{m}\fallf{x-m}^{n}.
\]

\section{Графика, таблицы и листинги}
\label{sec:floats}

\begin{figure}[htbp]
  \centering
  \begin{tikzpicture}[>=stealth, node distance=30mm]
    \node[draw, circle, fill=blue!15] (a) {$0$};
    \node[draw, circle, fill=green!15, right=of a] (b) {$1$};
    \node[draw, circle, fill=orange!15, below=of b] (c) {$2$};
    \draw[->, thick] (a) to[bend left] node[above] {$p$} (b);
    \draw[->, thick] (b) to[bend left] node[below] {$q$} (a);
    \draw[->, thick] (b) -- node[right] {$1-q$} (c);
    \draw[->, thick] (c) edge[loop right] node {$1$} (c);
  \end{tikzpicture}
  \caption{Небольшая цепь Маркова, набранная через TikZ из \texttt{flsuite}.}
  \label{fig:markov-chain}
\end{figure}

\begin{figure}[htbp]
  \centering
  \begin{tikzcd}[row sep=large, column sep=large]
    V \arrow[r] \arrow[d] & W \arrow[d] \\
    V' \arrow[r] & W'
  \end{tikzcd}
  \caption{Коммутативная диаграмма через \texttt{tikz-cd}.}
  \label{fig:diagram}
\end{figure}

\begin{figure}[htbp]
  \centering
  \begin{tikzpicture}
    \begin{axis}[
      width=0.78\linewidth,
      height=6cm,
      grid=both,
      xlabel={$x$},
      ylabel={$f(x)$},
      legend pos=north west
    ]
      \addplot[deepblue, thick, domain=-2:2, samples=100] {exp(-x^2)};
      \addplot[red!70!black, thick, domain=-2:2, samples=100] {x^2/4};
      \legend{$e^{-x^2}$, $x^2/4$}
    \end{axis}
  \end{tikzpicture}
  \caption{График, построенный средствами PGFPlots.}
  \label{fig:plot}
\end{figure}

\begin{table}[htbp]
  \centering
  \caption{Несколько операторов из \texttt{xamsmath}.}
  \label{tab:operators}
  \begin{tabular}{lll}
    \toprule
    Область & Команда & Результат \\
    \midrule
    Алгебра & \verb|\rk(A)| & $\rk(A)$ \\
    Анализ & \verb|\grad f| & $\grad f$ \\
    Вероятность & \verb|\D[X]| & $\D[X]$ \\
    Комбинаторика & \verb|\multbinom{n}{k}| & $\multbinom{n}{k}$ \\
    \bottomrule
  \end{tabular}
\end{table}

\begin{listing}[htbp]
\begin{minted}{python}
def total_probability(priors, conditionals):
    """Return sum_i P(A | H_i) P(H_i)."""
    return sum(p * c for p, c in zip(priors, conditionals))

priors = [0.2, 0.5, 0.3]
conditionals = [0.1, 0.4, 0.8]
print(total_probability(priors, conditionals))
\end{minted}
\caption{Минимальный расчет формулы полной вероятности.}
\label{lst:recursion}
\end{listing}

\section{refsuite: ссылки, библиография и индексы}
\label{sec:references}

Конфигурация \texttt{refsuite} в преамбуле этого файла написана через l3-ключи.
Она создает короткие команды ссылок, подключает \texttt{biblatex}, включает
\texttt{zref-clever} и настраивает два указателя. Классические источники по
\TeX{} и \LaTeX{} удобно цитировать прямо в лекции~\cite{knuth1984texbook,lamport1994latex},
а графические примеры естественно связываются с TikZ~\cite{tantau2013tikz,ctanpgf}.

\index{refsuite}\index{короткие ссылки}\index[notation]{$\P(A \mid B)$}
Ссылки читаются как обычный текст: \picref{fig:plot} содержит график,
\tabref{tab:operators} перечисляет команды, а \lstref{lst:recursion} показывает
короткий вычислительный фрагмент.

\clearpage
\chapter*{Справочные списки}
\addcontentsline{toc}{chapter}{Справочные списки}

\listofkeytheorems[title={Список теорем}]

\printbibliographymain
\printbibliographyonline

\printindexnotation
\printindex

\end{document}
